% L35 P3-OBS Pacer — printable hard copy of the live web pacer.
% Mirrors the admin's Lesson Plan Template (Teachers will / Students will /
% Questions to consider) with absolute time stamps for Period F (7:45--8:50).
%
% Compile: pdflatex --miktex-enable-installer L35_P3_obs_pacer.tex
\documentclass[10pt]{article}
\usepackage{preamble}

% Tighter margins for the pacer (more content per page)
\geometry{letterpaper, top=0.55in, bottom=0.7in, left=0.55in, right=0.55in}

\usepackage{fancyhdr}
\pagestyle{fancy}
\fancyhf{}
\renewcommand{\headrulewidth}{0.4pt}
\fancyhead[L]{\textbf{L35 P3-OBS Pacer}}
\fancyhead[R]{\textit{Fri 5/8/26 \textbullet{} Period F \textbullet{} 7:45--8:50}}
\fancyfoot[C]{\small Page \thepage{} of \pageref{LastPage} \quad \textbullet{} \quad Single source of truth: \texttt{web/pacers/L35\_Pacer.html} (P3-OBS tab)}
\usepackage{lastpage}

\setlength{\parskip}{4pt}

% ── Phase block macro ──────────────────────────────────────────────────────
% \phasehead{N}{title}{time range}{DOK}{minutes}{framework descriptor}
\newcommand{\phasehead}[6]{%
  \vspace{0.25em}
  \noindent\colorbox{tealheader}{\parbox{\dimexpr\linewidth-2\fboxsep}{%
    \color{white}\bfseries #1.\ #2 \hfill \normalfont\bfseries #3 \textbullet{} DOK #4 \textbullet{} #5 min%
  }}\par
  {\small\itshape\color{framegray}#6}\par
  \smallskip
}

% \itembox{label}{body}
\newcommand{\itembox}[2]{%
  \begin{tcolorbox}[
    enhanced, breakable, colback=frameshade, colframe=dayblue,
    boxrule=0.6pt, arc=2pt, left=7pt, right=7pt, top=4pt, bottom=4pt
  ]
  \textbf{Item:} \textbf{#1}\par#2
  \end{tcolorbox}
}

\begin{document}

% ── Banner ──────────────────────────────────────────────────────────────────
\daybanner%
{Observation Lesson \textbullet{} 65 min \textbullet{} Mirrors admin's Lesson Plan Template (Teachers will / Students will / Questions to consider)}%
{Lesson 3-5 \textbullet{} Day 3 \textbullet{} Apply / Performance Task}

\begin{tcolorbox}[
  enhanced, colback=calloutyellow, colframe=warmred,
  boxrule=0.5pt, arc=2pt, left=8pt, right=8pt, top=4pt, bottom=4pt
]
\textbf{Phase plan (65 min):} 7:45--7:55 Do Now (10) \textbullet{} 7:55--8:00 Launch (5) \textbullet{} 8:00--8:15 Fireworks (15) \textbullet{} 8:15--8:35 Venetta Performance Task (20) \textbullet{} 8:35--8:45 Share/Summary (10) \textbullet{} 8:45--8:50 Exit (5)
\end{tcolorbox}

\begin{tcolorbox}[
  enhanced, colback=lightgray, colframe=framegray!50,
  boxrule=0.5pt, arc=2pt, left=8pt, right=8pt, top=4pt, bottom=4pt
]
\textbf{Math Obj:} (1) Identify zeros via factoring or synthetic division. (2) Use zeros to graph a polynomial. (3) Interpret features in a real-world context.\par
\textbf{Language Obj:} Explain solutions verbally and in writing using sentence frames + structured talk-time.\par
\textbf{Essential Q:} How are the zeros of a polynomial related to its equation and graph?
\end{tcolorbox}

% ════════════════════════════════════════════════════════════════════════════
% PHASE 1 — DO NOW
% ════════════════════════════════════════════════════════════════════════════
\phasehead{1}{Do Now --- SAT practice from 3-5}{7:45--7:55}{2}{10}{Less than 10 min \textbullet{} low-floor entry}

\itembox{SAT practice (3-5)}{Without using a graphing calculator, determine which of the given graphs matches $f(x) = x^3 + x^2 - 4x$. Explain why you made that choice.}

\begin{callout}{\IconSpeak{} Teachers will\ldots}{calloutblue}
\begin{itemize}[leftmargin=1.4em,itemsep=1pt,topsep=1pt]
  \item Display problem on board
  \item Read prompt
  \item Tell students they have 3-5 minutes to work on problem (*no calculator)
  \item Set timer for 3 minutes independently think/write
  \item After timer, 3 minutes to pair up and share answers with each other
  \item After 5 minutes, class discussion
\end{itemize}
\end{callout}

\begin{callout}{[STUDENT] Students will\ldots}{calloutell}
\begin{itemize}[leftmargin=1.4em,itemsep=1pt,topsep=1pt]
  \item Actively listen to prompt
  \item Record initial ideas on their paper (3 min)
  \item Share with a partner and record explanation for their choice (3 min)
  \item Share out in class discussion (2-3)
\end{itemize}
\end{callout}

\begin{callout}{\IconPuzzle{} Questions to consider\ldots}{calloutyellow}
\begin{itemize}[leftmargin=1.4em,itemsep=1pt,topsep=1pt]
  \item How can you engage students in the lesson?
  \item How can you hook them?
  \item How are all adults intentionally planned for?
\end{itemize}
\end{callout}

\begin{callout}{\IconStar{} Answer key}{calloutgreen}
Factor: $f(x) = x(x^2 + x - 4)$. Zero at $x=0$ from the linear factor; quadratic factor has discriminant $1 - 4(1)(-4) = 17 > 0$, giving $x \approx -2.56$ and $x \approx 1.56$. Three real zeros total. Degree 3 + positive leading coefficient $\Rightarrow$ end behavior DOWN--UP. Eliminate graphs that miss $x=0$, show wrong end behavior, or have fewer than 3 zeros.
\end{callout}

% ════════════════════════════════════════════════════════════════════════════
% PHASE 2 — LAUNCH
% ════════════════════════════════════════════════════════════════════════════
\phasehead{2}{Launch --- Unpack objectives + connect to prior}{7:55--8:00}{1-2}{5}{5--15 min \textbullet{} academic conversation}

\begin{callout}{\IconSpeak{} Teachers will\ldots}{calloutblue}
\begin{itemize}[leftmargin=1.4em,itemsep=1pt,topsep=1pt]
  \item Teacher asks for volunteers to read objectives
  \item Connect to previously learned skills / why we're doing what we're doing
\end{itemize}
\end{callout}

\begin{callout}{[STUDENT] Students will\ldots}{calloutell}
\begin{itemize}[leftmargin=1.4em,itemsep=1pt,topsep=1pt]
  \item Volunteers read objectives
\end{itemize}
\end{callout}

\begin{callout}{\IconPuzzle{} Questions to consider\ldots}{calloutyellow}
\begin{itemize}[leftmargin=1.4em,itemsep=1pt,topsep=1pt]
  \item What will you explicitly model and how?
  \item How will you know if students are ready to work alone?
  \item How are all adults intentionally planned for?
\end{itemize}
\end{callout}

% ════════════════════════════════════════════════════════════════════════════
% PHASE 3 — EXPLORE: FIREWORKS
% ════════════════════════════════════════════════════════════════════════════
\phasehead{3}{Explore --- Fireworks Problem}{8:00--8:15}{2}{15}{15 min \textbullet{} TWPS \textbullet{} cognitively demanding}

\itembox{Fireworks: $h(t) = -16t^2 + 32t + 6.5$}{Three parts: (a) reasonable domain, (b) zeros + meaning, (c) vertex + meaning. Apply / Make Sense and Persevere section of the student packet.}

\begin{callout}{\IconSpeak{} Teachers will\ldots}{calloutblue}
\begin{itemize}[leftmargin=1.4em,itemsep=1pt,topsep=1pt]
  \item Teacher asks student to read the problem(s). Cue up students to TWPS.
  \item \textbf{Questions to ask:} Discuss domain. Why is it important to think about the domain before starting? Why would only values $> 0$ work? What is domain talking about in this case? What are possible values for $t$? Can we have negative time? Why/why not?
  \item Instruct students 3-4 minutes to think/pair and prepare to share with class
  \item Circulate while students are thinking/pairing to identify students to warm-call to share out
  \item Ask students to share out
  \item Display on Desmos
  \item Release students to work on B and C in pairs. Communicate clear expectations (5-min timer to let them productively struggle). Add 2-3 more minutes as needed
  \item Circulate to ask open questions, assess progress, redirect as needed, clarify common misconceptions as a class
  \item Ask students to share their thinking (whoever did not participate in part A)
\end{itemize}
\end{callout}

\begin{callout}{[STUDENT] Students will\ldots}{calloutell}
\begin{itemize}[leftmargin=1.4em,itemsep=1pt,topsep=1pt]
  \item Actively listen to problem being read
  \item Respond to questions the teacher asks
  \item Think independently (1 min)
  \item Pair up w/partner and share thinking (2 min)
  \item Students share their thinking when called on (2-3 min)
\end{itemize}
\end{callout}

\begin{callout}{\IconStar{} Answer key}{calloutgreen}
\textbf{(a) Domain:} $0 \le t \le \approx 2.19$ s (launch to ground impact). $t < 0$ has no physical meaning.\par
\textbf{(b) Zeros:} Solve $-16t^2 + 32t + 6.5 = 0 \Rightarrow t \approx 2.19$ s (the negative root $\approx -0.19$ s is rejected). The positive zero is the moment the firework hits the ground.\par
\textbf{(c) Vertex:} $t = -b/(2a) = -32/(-32) = 1$ s. $h(1) = -16 + 32 + 6.5 = 22.5$ m. Vertex represents \textbf{maximum height} (22.5 m at $t = 1$ s).
\end{callout}

% ════════════════════════════════════════════════════════════════════════════
% PHASE 4 — EXPLORE: VENETTA PERFORMANCE TASK (DOK-3 DRIVER)
% ════════════════════════════════════════════════════════════════════════════
\phasehead{4}{Explore --- Performance Task: Venetta Deli {[STAR] DOK-3 DRIVER}}{8:15--8:35}{3}{20}{15--20 min \textbullet{} TPS \textbullet{} DOK-3 driver \textbullet{} *Connect back to PBL}

\itembox{Venetta profit: $P(t) = t(t+3)(t-2)$ \quad {\IconStar} DOK-3}{Profit in hundreds of dollars, $t$ years after 2000. \textbf{(A)} Sketch + confirm with Desmos. \textbf{(B)} During what years did Venetta NOT make a profit? \textbf{(C)} Predict 2020 profit if the model is appropriate.}

\begin{callout}{\IconSpeak{} Teachers will\ldots}{calloutblue}
\begin{itemize}[leftmargin=1.4em,itemsep=1pt,topsep=1pt]
  \item \textbf{Connect it to real-world:} ``Does anyone know anyone that owns or operates any food businesses/restaurants/shops?''
  \item ``Well, it can be stressful---business owners aren't always profitable. Why could that be? This is what this could look like in real life---but with Algebra.''
  \item ``Does anyone remember the profit equation from PBL/project?'' *Connect back to PBL
  \item Teacher frames the problem (3-4 min): here's another profit equation for another store (not pizza, but deli sandwiches this time)
  \item What variables do we see here? What does $t$ mean?
  \item Read task itself and then go to part A (5 min). \textbf{Question:} How can you identify the zeros from the factored form? How can you determine the behavior?
  \item \textbf{Summarize:} show student work from iPad and ask students ``what do you think about this?'' ``do you agree or disagree?''
  \item B/C---Read and release students to think/pair/share their thinking ($\sim$7 min)
  \item \textbf{Question:} How do you know, looking at a graph, if you have a profit?
  \item For part C, explore a variety of pathways to get that answer
\end{itemize}
\end{callout}

\begin{callout}{[STUDENT] Students will\ldots}{calloutell}
\begin{itemize}[leftmargin=1.4em,itemsep=1pt,topsep=1pt]
  \item Engage in discussion around food businesses/profits/etc.
  \item Recall prior knowledge to PBL project
  \item Work in pairs/triads on B and C
  \item Engage in discussion around B and C with one another, and respond to questions from teacher as needed
  \item Share their thinking with the class
\end{itemize}
\end{callout}

\begin{callout}{\IconPuzzle{} Questions to consider\ldots}{calloutyellow}
\begin{itemize}[leftmargin=1.4em,itemsep=1pt,topsep=1pt]
  \item What will students do for most of the session? What will they practice? How?
  \item What type of academic conversations will students engage in?
  \item How will you be interacting with them?
\end{itemize}
\end{callout}

\begin{callout}{\IconStar{} Answer key}{calloutgreen}
\textbf{(A) Zeros:} $t = 0, -3, 2$. Only $t \ge 0$ meaningful. Graph dips below $t$-axis on $0 < t < 2$, crosses up at $t = 2$, climbs as a cubic with $+$ leading coefficient.\par
\textbf{(B) No profit:} on $0 < t < 2 \Rightarrow$ \textbf{years 2000--2001} (broke even at start of 2002).\par
\textbf{(C) 2020 prediction:} $t = 20 \Rightarrow P(20) = 20 \cdot 23 \cdot 18 = 8280$ hundreds of dollars $= \$828{,}000$. \textbf{CAVEAT (DOK-3 pivot):} 20-year extrapolation overshoots dramatically---cubic growth, market saturation, competition. \emph{Discuss model limits.}
\end{callout}

\begin{callout}{\IconWarn{} Warning / Variant}{calloutred}
\textbf{Part C is the DOK-3 pivot} --- students must JUDGE model validity, not just compute $P(20)$. Press the model-limits discussion or this stays at DOK-2.
\end{callout}

% ════════════════════════════════════════════════════════════════════════════
% PHASE 5 — SHARE / SUMMARY
% ════════════════════════════════════════════════════════════════════════════
\phasehead{5}{Share / Summary}{8:35--8:45}{2}{10}{5--15 min \textbullet{} synthesis \textbullet{} return to objectives}

\begin{callout}{\IconSpeak{} Teachers will\ldots}{calloutblue}
\begin{itemize}[leftmargin=1.4em,itemsep=1pt,topsep=1pt]
  \item After $\sim$7 min, call group back together to report their thinking/findings. \textbf{Ensure equity of voice.}
  \item Share key takeaways
\end{itemize}
\end{callout}

\begin{callout}{[STUDENT] Students will\ldots}{calloutell}
\begin{itemize}[leftmargin=1.4em,itemsep=1pt,topsep=1pt]
  \item Report thinking and findings to class
  \item Engage in equity of voice --- listen and let others share
\end{itemize}
\end{callout}

\begin{callout}{\IconPuzzle{} Questions to consider\ldots}{calloutyellow}
\begin{itemize}[leftmargin=1.4em,itemsep=1pt,topsep=1pt]
  \item How will students express what they have learned? This should be directly related to the mini-lesson and objective.
\end{itemize}
\end{callout}

\begin{callout}{\IconChain{} Bridge}{calloutpeach}
Bridge to Mon 5/11 (P3-cont, Storage Box \#27): same factored-form reasoning, but now zeros are physical dimensions of a box.
\end{callout}

% ════════════════════════════════════════════════════════════════════════════
% PHASE 6 — EXIT TICKET
% ════════════════════════════════════════════════════════════════════════════
\phasehead{6}{Exit Ticket --- Baseball + reflection}{8:45--8:50}{2}{5}{5 min \textbullet{} DOK-2 cool-down}

\itembox{Baseball --- explain why graph shows only one zero}{Height of a baseball thrown in the air, modeled by a quadratic. \textbf{Explain why the graph shows only ONE zero.} Hint: domain restriction.}

\begin{callout}{\IconSpeak{} Teachers will\ldots}{calloutblue}
\begin{itemize}[leftmargin=1.4em,itemsep=1pt,topsep=1pt]
  \item Project exit prompt. Students complete independently and turn in at the door
  \item Tonight: scan answers; if many miss the domain-restriction idea, open Mon 5/11 with a Do Now bridge
\end{itemize}
\end{callout}

\begin{callout}{[STUDENT] Students will\ldots}{calloutell}
\begin{itemize}[leftmargin=1.4em,itemsep=1pt,topsep=1pt]
  \item Complete the Baseball reflection independently
  \item Submit at the door
\end{itemize}
\end{callout}

\begin{callout}{\IconPuzzle{} Questions to consider\ldots}{calloutyellow}
\begin{itemize}[leftmargin=1.4em,itemsep=1pt,topsep=1pt]
  \item How will students show what they've learned based on the objective?
\end{itemize}
\end{callout}

\begin{callout}{\IconStar{} Answer key}{calloutgreen}
A quadratic can have two mathematical zeros, but only one is in the meaningful domain ($t \ge 0$). The other zero is at \textbf{negative time} --- before the ball was thrown --- which has no physical meaning. The graph shows only the zero where the ball lands. Same idea as Fireworks (a) and Venetta (B): real-world contexts restrict which mathematical zeros are physically relevant.
\end{callout}

\end{document}
